We collect 75 availability issues from 28 subsystems of the e-commerce system of Alibaba.
These subsystems contain 265 services on average (min 7, max 1687, median 82).
The time of these availability issues ranges from February to June 2020.
All of them have been identified and annotated with anomaly types and root causes by operation engineers.
Each of them may have multiple anomaly types, and each anomaly type has a single root cause.
Among the 75 availability issues, 37 are caused by performance anomaly, 43 are caused by reliability anomaly, and 21 are caused by traffic anomaly.
%Note that an availability issue may be caused by multiple types of anomaly problems.
The construction of the service call graph follows the practice in Alibaba, i.e., collecting service calls in the last 30 minutes and service call metrics (i.e., RT, EC, and QPS) in the last 60 minutes.

%of the time window for extracting the calling relationship between services of the microservice system and extracting metrics data are half an hour and one hour before the alarm occurs. In the scenario of this article, these two time windows can well cover the service nodes with calling relationship and the abnormal metrics data when the alarm occurs.

To evaluate the accuracy and efficiency, we compare \app~with the following two state-of-the-art baseline approaches.

\begin{itemize}
\item \textbf{MonitorRank}~\cite{cscs2013monitorRank}:
It detects root causes of anomalies in service-oriented systems and provides a ranked order list of possible root causes.
It uses the historical and current time-series metrics of each service along with service call graph to build an unsupervised model for ranking.

%It builds the service call graph through the Hadoop tools.
%To diagnose root cause services, it firstly calculates the correlation between the front-end service which reported with alarm and each node service in the graph as an abnormal similarity score. Then it records the visiting probability of each node during the random walk process in the graph. Finally, the abnormal nodes are ranked with the customized PageRank algorithm and output the top k items. In order to reproduce MonitorRank, we replicate the random walk and customized PageRank algorithm, and the abnormal similarity score is calculated with Pearson correlation coefficient of the metrics of alarm service and each node in the call graph.

\item \textbf{Microscope}~\cite{ICSOC2018Microscope}:
It locates anomalous services with a ranked list of possible root causes in microservice systems.
It builds and uses service instance causality graph, which captures both communicating (service calls) and non-communicating (services co-located in the same machines) service dependencies.
It traverses the graph from the front-end service where an anomaly is reported and ranks possible root causes based the similarity of their metrics with those of the front-end service.
%The metrics used by Microscope is collected from microservice instance, while the metrics that we obtain from infrastructures in Alibaba is aggregated to the microservice level. It is hard to distinguish which microservice instances are deployed on the same machine in the huge industrial-scale microservice system. So we use the metrics of microservice level to build the service dependency graph and detect abnormal microservices.
\end{itemize}

%Note that Microscope is targeted at performance anomaly but may also be used for
% we think the its approach is general and can be applied to detect the other two anomalies. So we also conduct the corresponding experiments with Microscope.

We use the top-$k$ hit ratio (HR@$k$) and mean reciprocal rank (MRR) to measure the accuracy of the root case localization.
\begin{itemize}

	\item HR@$k$ denotes the probability that the top $k$ result list contains the root cause.
	Let $A$ be the set of availability issues, $r_i$ be the root cause of the $i$th availability issue, $Rank_{i}^{k}$ be the top $k$ result list of the $i$th availability issue, HR@$k$ can be calculated as follows.

	\begin{equation}
	HR@k = \cfrac{1}{|A|}\sum\limits_{i=1}^{|A|}(r_i \in Rank_{i}^{k})
	\end{equation}

	\item MRR is the multiplicative inverse of the rank of the first correct answer.	
	If the correct answer is not included in the returned list, the rank can be regarded as positive infinity.
	Let $A$ be the set of availability issues and $Index_i$ be the rank of the root cause in the returned list of the $i$th availability issue, MRR can be calculated as follows.
	
	\begin{equation}
	MRR = \cfrac{1}{|A|}\sum\limits_{i=1}^{|A|}\cfrac{1}{Index_i}
	\end{equation}
	
\end{itemize}

All the experiments are conducted on a cluster with 15 virtual machines on Alibaba Cloud.
Each virtual machine is equipped with 4 Intel Xeon 2.50GHz CPU, 8 GB RAM and 60 GB Disk, and running Alibaba Group Enterprise Linux Server release 6.2.
The default settings of the thresholds are the following:
the RT threshold used in reliability anomaly detection is 50ms;
the correlation threshold used in the traffic anomaly detection is 0.9;
the correlation threshold used in the pruning strategy is 0.7.
